\documentclass[12pt,a4paper]{article}
\usepackage[utf8]{inputenc}
\usepackage[vietnamese]{babel}
\usepackage{amsmath,amssymb}
\usepackage[margin=2cm]{geometry}
\usepackage{enumitem}
\usepackage{fontspec}
\setmainfont{Times New Roman}

\title{\textbf{ĐÁP ÁN VÀ LỜI GIẢI CHI TIẾT}\\
\large Bài 4 -- Nhiệt Dung Riêng -- GEMS Vật Lý 12}
\author{Giáo viên: Kha Khung Hiệp}
\date{}

\begin{document}
\maketitle

\section*{PHẦN I: TRẮC NGHIỆM}
\begin{enumerate}
  \item \textbf{C} (Theo định nghĩa hệ SI: đơn vị nhiệt dung riêng là J/kg.K).
  \item \textbf{B} (Nhiệt lượng để 1 kg nhôm tăng thêm 1 K là 880 J).
\end{enumerate}

\section*{PHẦN II: ĐÚNG/SAI}
\textbf{Câu 1:}
\begin{enumerate}[label=\alph*)]
  \item \textbf{ĐÚNG} (Vỏ nhựa cách nhiệt tốt giúp giảm nhiệt truyền ra bên ngoài).
  \item \textbf{SAI} (Công thức đúng là $c = \dfrac{P \cdot t}{m \cdot \Delta t}$, không phải $\dfrac{P \cdot m}{t \cdot \Delta t}$).
  \item \textbf{SAI} (Nhiệt thất thoát ra ngoài làm nhiệt độ tăng chậm hơn, $\Delta t$ nhỏ hơn nên giá trị $c$ tính ra \textit{lớn hơn} giá trị thực, không phải nhỏ hơn).
  \item \textbf{ĐÚNG} (Khuấy đều giúp phân bố nhiệt độ đồng đều, giảm thiểu sai số đo).
\end{enumerate}

\section*{PHẦN III: CÂU HỎI TRẢ LỜI NGẮN}
\textbf{Câu 1 -- Lời giải:}
\begin{itemize}
  \item $\Delta t = 85 - 25 = 60^{\circ}\text{C}$ (hoặc 60 K).
  \item $c = \dfrac{Q}{m \cdot \Delta t} = \dfrac{11400}{0{,}5 \cdot 60} = \dfrac{11400}{30} = 380\text{ J/kg.K}$.
  \item \textbf{Đáp số: 380}
\end{itemize}

\vfill
\begin{center}
\textit{--- GEMS Vật Lý 12 -- Bài 4: Nhiệt Dung Riêng ---}
\end{center}

\end{document}
